\documentclass[11pt]{article}
\usepackage[T1]{fontenc}
\usepackage[utf8]{inputenc}
\usepackage[english]{babel}
\usepackage{lmodern,amsmath,amssymb,amsthm,booktabs,microtype}
\usepackage[margin=1in]{geometry}
\usepackage[colorlinks=true,linkcolor=black,citecolor=blue,urlcolor=blue]{hyperref}
\hypersetup{pdftitle={Permutations from Fourth-Order Binomial Window Sums},
  pdfauthor={Carlo Corti},pdfsubject={Exact enumeration of OEIS A180215 through n=30}}
\pdfinfoomitdate=1
\pdftrailerid{}
\pdfsuppressptexinfo=15
\newtheorem{theorem}{Theorem}[section]
\newtheorem{proposition}[theorem]{Proposition}
\newcommand{\NN}{\mathbb Z_{\geq 0}}
\newcommand{\rev}{\operatorname{rev}}
\title{Permutations from Fourth-Order Binomial Window Sums:\\[2pt]
OEIS A180215 through \(n=30\)}
\author{Carlo Corti\\Independent researcher\\\texttt{carlo.corti@outlook.com}}
\date{November 2026}

\begin{document}
\maketitle

\begin{abstract}
We count permutations of \(0,\ldots,n-1\) realized by consecutive
weighted sums of nonnegative integers with coefficients \(1,4,6,4,1\).
The output value zero forces five consecutive input entries to vanish.
The rest of the input is uniquely determined, and the permutation splits
into two outward words on complementary sets of labels. We derive an exact
count from this decomposition and give two necessary constraints on the
four-coordinate state of a word. Every fourth-order realizable permutation
is also realizable for the third- and second-order kernels. A depth-first
enumeration and a separate weighted-state dynamic program using the same
decomposition agree coefficient by coefficient and give
all counts through \(n=30\), extending Hardin's 24-term baseline by six
terms, with \(a(30)=3209944\). A direct enumeration of input vectors checks
the original definition separately through \(n=24\).
\end{abstract}

\noindent\textbf{2020 Mathematics Subject Classification.}
Primary 05A05. Secondary 05A15, 05A19, 11Y55.

\smallskip
\noindent\textbf{Key words and phrases.}
Permutations, binomial weights, nonnegative integer realizations,
subset enumeration, dynamic programming, OEIS A180215.

\section{The inverse problem}

For \(n\geq 1\), let \(p=(p_0,\ldots,p_{n-1})\) be a permutation of
\(\{0,\ldots,n-1\}\). A \emph{realization} of \(p\) is a vector
\(x=(x_0,\ldots,x_{n+3})\in\NN^{n+4}\) satisfying
\begin{equation}
 p_i=x_i+4x_{i+1}+6x_{i+2}+4x_{i+3}+x_{i+4},
 \qquad 0\leq i<n.
 \label{eq:definition}
\end{equation}
Write \(a(n)\) for the number of permutations admitting a realization.
A permutation and its reversal are counted separately whenever they
differ. The entries of \(x\) may repeat.

Hardin introduced this sequence as OEIS A180215 on 15 August 2010
\cite{hardin215}. The historical data available at the start of this
study comprised the 24 values ending with \(a(24)=12944\). These values
appear in revision~1 of the entry history. The OEIS data mirror and the
b-file linked from the entry, checked on 21 November 2026, give the
same prefix \cite{hardin215}. Hardin's adjacent entries A180213 and A180214 give
the analogous problems with kernels \(1,2,1\) and \(1,3,3,1\)
\cite{hardin213,hardin214}. We derive a counting method for the
fourth-order problem, compare the three sets of realizable permutations
at fixed length, and extend the 24-term baseline through \(n=30\).

\section{Cutting at the zero window}

If \(p_\ell=0\), positivity of the coefficients and nonnegativity of \(x\)
force
\[
 x_\ell=x_{\ell+1}=x_{\ell+2}=x_{\ell+3}=x_{\ell+4}=0.
\]
Starting from that window and moving right, each next equation determines
one new rightmost entry, whose coefficient is \(1\). Moving left,
each equation determines one new leftmost entry, also with coefficient
\(1\). These recurrences reach both ends of the input even when
\(\ell=0\) or \(\ell=n-1\), in which case one sweep is empty. Thus any
realizable permutation has exactly one realization.

An \emph{admissible word} \(v=(v_1,\ldots,v_k)\) consists of distinct
positive labels for which the recurrence
\begin{equation}
 v_j=y_j+4y_{j-1}+6y_{j-2}+4y_{j-3}+y_{j-4},
 \qquad 1\leq j\leq k
 \label{eq:side}
\end{equation}
gives \(y_j\geq0\) at every step, including \(j=k\), starting with
\(y_{-3}=y_{-2}=y_{-1}=y_0=0\). There is no extension condition after
the last label. For a finite set \(S\) of positive integers, let \(F(S)\) be the
number of admissible orderings of \(S\). The empty word gives
\(F(\varnothing)=1\). The recurrence uses actual integer labels. For
instance, \((4,16)\) is admissible with \((y_1,y_2)=(4,0)\), whereas
\((1,2)\), its relative-order standardization, fails the second step.

\begin{theorem}\label{thm:split}
With \(U_m=\{1,\ldots,m\}\) and \(U_0=\varnothing\), we have
\begin{equation}
 a(n)=\sum_{S\subseteq U_{n-1}}F(S)\,F(U_{n-1}\setminus S)
 \qquad (n\geq1).
 \label{eq:sum}
\end{equation}
The subset \(S\) is the set of labels to the left of zero.
\end{theorem}

\begin{proof}
Read the labels outward on each side of zero, taking
\(v^{\mathrm L}=(p_{\ell-1},\ldots,p_0)\) and
\(v^{\mathrm R}=(p_{\ell+1},\ldots,p_{n-1})\).
For \(j\geq1\) on the corresponding side, put
\(y^{\mathrm L}_j=x_{\ell-j}\) and
\(y^{\mathrm R}_j=x_{\ell+4+j}\).
For either side the four padded entries lie in the zero block.
At \(i=\ell-j\) the forward window uses the left coordinates from
\(y^{\mathrm L}_j\) down to \(y^{\mathrm L}_{j-4}\). At \(i=\ell+j\)
it uses the right coordinates from \(y^{\mathrm R}_{j-4}\) up to
\(y^{\mathrm R}_j\). Palindromicity gives \eqref{eq:side} in both cases.

Conversely, choose a left word on \(S\) and a right word on
\(U_{n-1}\setminus S\). With \(\ell=|S|\), place five zeros at positions
\(\ell,\ldots,\ell+4\), put the left realization entries in reverse
order before them, and put the right entries in outward order after them.
The vector has \(\ell+5+(n-1-\ell)=n+4\) entries. The window at
\(\ell\) has value zero, and its neighbours, when present, use one
side entry and four zeros. In general the windows at \(\ell\pm j\)
give the two recurrences just derived. The output is
\((\rev(v^{\mathrm L}),0,v^{\mathrm R})\), with the words concatenated
around zero. Cutting and gluing are inverse also for either empty word.
For \(n=1\) the two words are empty and the unique realization is the
five-entry zero vector, giving \(a(1)=1\). For fixed \(S\), the position
of zero and the left and right choices give exactly the product in
\eqref{eq:sum}.
\end{proof}

For \(n=3\), the words \((1)\) and \((2)\) give
\(p=(1,0,2)\) with \(x=(1,0,0,0,0,0,2)\). Swapping the words gives
its reversal. A second label on either side would have to be at least
four times the first by \eqref{eq:side}. Thus if \(n\le4\), each side
has at most one label. In particular, at \(n=4\) the two sides cannot
hold all three positive labels, so \(a(4)=0\).

Equation \eqref{eq:sum} evaluates the subset convolution of \(F\) with
itself at the nested sets \(U_0,U_1,\ldots\).
Subset convolution is commonly defined with one output for each subset
\cite[Eq.~(1)]{subsetconv}. Here the two implementations evaluate only the
required nested outputs, by direct complementary summation.

\begin{proposition}\label{prop:reversal}
For \(n>1\), reversal is a fixed-point-free involution on the realizable
permutations. In particular, \(a(n)\) is even.
\end{proposition}

\begin{proof}
Set \(x'_j=x_{n+3-j}\) for \(0\le j\le n+3\). Its window at \(i\)
reads the original window at \(n-1-i\) in reverse order, with the same
value \(p_{n-1-i}\) by palindromicity. Applying reversal twice recovers
the original permutation. A fixed point with \(n>1\) would have equal
first and last labels.
\end{proof}

For \(m\geq1\), \(S\) and \(U_m\setminus S\) give equal products in
\eqref{eq:sum}. Exactly one subset in each pair omits \(m\).
Doubling the terms indexed by those subsets accounts for both ordered
choices of sides.

\section{Four-coordinate states and exact multiplicities}

Fix a maximum label \(M\) and use the universe \(U_M\).
After a word of length \(k\), store its used set \(S\) and the four latest
coordinates in chronological order,
\((\alpha,\beta,\gamma,\delta)=(y_{k-3},y_{k-2},y_{k-1},y_k)\), with the same zero padding
as in \eqref{eq:side}. The starting state is
\((\varnothing,0,0,0,0)\). Define \(B=\alpha+4\beta+6\gamma+4\delta\).
An unused \(v\in U_M\) can extend the word exactly when \(v\ge B\).
Then the new coordinate is \(e=v-B\), and the new state is
\begin{equation}
 (S\cup\{v\},\beta,\gamma,\delta,e).
 \label{eq:transition}
\end{equation}
The next word label is subject to \(B\), but the current word is already
valid. For example, the word \((1)\) contributes to \(F(\{1\})\)
when \(M=1\), even though the next-step bound is \(4\).

\begin{proposition}\label{prop:bound}
For a reachable state after \(k\) labels from \(U_M\),
\begin{align}
 4\alpha+6\beta+4\gamma+\delta &\le M, \label{eq:last}\\
 \sum_{v\in S}v-(15\alpha+11\beta+5\gamma+\delta)
 &\in 16\NN. \label{eq:mass}
\end{align}
\end{proposition}

\begin{proof}
For \(k>0\), the last label is
\(y_{k-4}+4\alpha+6\beta+4\gamma+\delta\), giving \eqref{eq:last}. In the sum
of the first \(k\) labels, each \(y_j\) with \(1\le j\le k-4\) has
received all five weights, with sum \(16\). The four boundary coordinates
have accumulated weights \(15,11,5,1\), respectively. Thus
\[
 \sum_{v\in S}v=16\sum_{j=1}^{k-4}y_j+15\alpha+11\beta+5\gamma+\delta.
\]
The sum over \(y_j\) is empty for \(k\le4\), and the padding gives both assertions
for the empty word.
\end{proof}

These conditions are necessary. For example, with \(M=16\), the hypothetical
state \((\{16\},0,0,0,0)\) satisfies both. The sole ordering on that
set reaches \((\{16\},0,0,0,16)\), so the hypothetical state is unreachable.

The dynamic program uses the exact transition \eqref{eq:transition}.
Write \(q\) for a state, \(D(q)\) for its word multiplicity, and
\(q\xrightarrow{v}q'\) for the legal extension by \(v\). The empty
state has multiplicity \(1\). At every positive cardinality,
\[
 D(q')=\sum_{(q,v):\,q\xrightarrow{v}q'}D(q).
\]
Deleting the last label of a word reaching \(q'\) recovers its
unique prefix and final transition. Conversely, each word reaching a
predecessor \(q\) can take the indicated transition. This proves the
recurrence by induction on \(|S|\), including when distinct words
merge into the same state. Summing over all states with set \(S\)
gives \(F(S)\), including states with no outgoing transition.

\section{Computing and checking the side counts}

\subsection{One depth-first enumeration for all prefixes}

The serial C17 producer traverses the tree of admissible words on
\(U_M\) using an explicit stack. Each node, including the empty word,
increments its mask cell on entry, before the next-label test. The
children are exactly the unused labels at least \(B\).
Every edge adds one label, so the tree has depth at most \(M\).
Truncating a word at its last label gives its unique parent, and the
transition recovers the word from that parent and label. These facts
establish exhaustiveness and a single visit per admissible word.

The table has \(2^M\) unsigned 128-bit cells. A coefficient depends
only on the labels of its set, so one table for \(M=29\) supplies
all terms \(1\leq n\leq30\). For a term of length \(n\), the prefix
size is \(m=n-1\). Bit \(v-1\) represents \(v\), and result-array
slot \(m\) holds \(a(m+1)\).
For \(n>1\) the producer sums products with the highest bit of
\(U_{n-1}\) absent on one side, then doubles the sum. There are
\(\sum_{m=1}^{M}2^{m-1}=2^M-1\) complementary pairs evaluated over
all those prefixes. The case \(n=1\) is initialized separately.

The dense storage and reduction counts are exponential in \(M\).
The number of word visits, including the root, is
\(W_M=\sum_{S\subseteq U_M}F(S)\). There are \(W_M-1\) tree edges.
The producer selects children by bit set. The elementary upper bound
\(\sum_{k=0}^M M!/(M-k)!\) includes sequences of distinct labels
regardless of reachability. Actual work at \(M=29\) is reported
in Table~\ref{tab:work}.
The practical reduction comes from starting at the forced zero window,
enumerating one-sided words and joining their counts instead of
constructing every pair, and restricting extensions to \(v\ge B\).
Word visits, subset cells and
final permutations count different objects.

\subsection{Separate implementation of the coefficient check}

The C++17 checker shares the zero-window theorem and the kernel with the
producer. It expands \eqref{eq:transition} by cardinality of \(S\).
It merges equal keys in a sparse hash table with additive multiplicities,
then sorts each completed frontier for comparison by mask. Each
coordinate is at most the label that produced it, so five bits per
coordinate suffice at the supported maxima \(n\le32\). With at most
31 label bits, each state key fits in 51 bits. For each cardinality,
the checker scans the dense table, sums all frontier weights at each
mask of that cardinality, and compares the result with the corresponding
producer cell. An absent mask has sum zero and is checked as such.
This covers every cell, including \(F(\varnothing)\).
The checker then uses the checked table to evaluate the full ordered sum
\eqref{eq:sum} for each \(n>1\). It compares these sums and the
initialized value \(a(1)=1\) with the producer output.

The programs differ in language, traversal, state representation, accumulation,
complementary reduction and persistence. The producer journals mask entries
and traversal cursors, whereas the checker persists complete frontiers during
coefficient verification and an ordered-sum cursor during reduction.
On loading a saved frontier, the checker tests \eqref{eq:last}
and \eqref{eq:mass} on each record. Transition generation uses the exact
criterion \(v\ge B\) and the unused-label test. The terminal checker
checkpoint has the final counts and counters and an empty frontier. It
contains no earlier frontiers. The check through \(n=30\) tests the
implementation of the common reduction. The traversal in the next
subsection bypasses that reduction through \(n=24\).

Unsigned 128-bit arithmetic is supplied by the GCC extension
\cite{gcc128}. The bound \(F(S)\le |S|!\) follows from the number of
orderings. Multiplicities are bounded by \(F(S)\), and nonnegative
products and partial sums in \eqref{eq:sum} are bounded by \(a(n)\le n!\).
In particular \(32!<2^{118}\), which bounds the counting arithmetic
through the source limit. A synthetic fixture with factorial rows exercised
wide-integer encoding and the output-finalization routines. On a second
call, the file-writing routines accepted the existing results, validation
report and completion record after comparing them byte for byte with the
expected output.
These are software tests on fabricated inputs. The sequence run had
maximum \(30\).

\subsection{Direct check from the original equations}

A separate C program enumerates feasible starting coordinates
\(x_0,\ldots,x_3\) with
\(x_0+4x_1+6x_2+4x_3\le n-1\). At each step it tries the nonnegative
choices for \(x_{i+4}\) that keep \(p_i\le n-1\), and discards repeated
output labels. If \(i+1<n\), a necessary condition for continuation is
\[
 x_{i+1}+4x_{i+2}+6x_{i+3}+4x_{i+4}\le n-1.
\]
The left side increases with \(x_{i+4}\), so larger choices can also be
discarded once the bound is exceeded. Every realization of
\eqref{eq:definition} survives these tests and appears as a leaf. A
leaf has \(n\) distinct labels in \(\{0,\ldots,n-1\}\), and uniqueness
of the realization makes the leaf count a permutation count.

This traversal starts from the original equations. The local
code-review integration run on 20 November 2026 reproduced Hardin's
counts for \(1\le n\le24\) with the bounded C program. The source,
build command and complete output are retained. Earlier Python
attempts at \(n=24\) were incomplete. For \(25\le n\le30\) the
verification is the coefficient comparison and ordered reduction
described above.

\section{Counts through \texorpdfstring{\(n=30\)}{n=30}}

Corti ran the producer and the checker on 20 November 2026, with
\(n=30\) as the chosen stopping point. The validated \(b\)-file
contains \(1\le n\le30\). The first uncomputed index in this study is
\(31\). The prefix through \(24\) agrees with Hardin's original
data \cite{hardin215}. Table~\ref{tab:extension} gives the six additional
values. The project materials are described in the availability statement
and cited in \cite{corti}.

\begin{table}[htbp]
\centering
\begin{tabular}{lrrrrrr}
\toprule
\(n\)    & 25 & 26 & 27 & 28 & 29 & 30 \\
\midrule
\(a(n)\) & 33572 & 50880 & 181968 & 464508 & 1150932 & 3209944 \\
\bottomrule
\end{tabular}
\caption{Exact extension of the historical 24-term prefix, with every
row checked by the separate coefficient census and ordered sum.}
\label{tab:extension}
\end{table}

Both phases used one thread on an AMD Ryzen 9 7940HS (x86-64).
The author reports Linux Mint 22.3 as the workstation distribution.
The retained host and build records document the processor, architecture,
GCC and G++ 13.3.0 with \texttt{-O3}, and OpenSSL 3.0.13 for
integrity digests.
The producer and checker took about 30 and 110 seconds, respectively.
The total was about 140 seconds of native wall time and 138 CPU seconds.
The two invocations, including their wrappers and post-run
integrity audits, took about 154 seconds in total.

\begin{table}[htbp]
\centering
\begin{tabular}{lr}
\toprule
Quantity at maximum \(n=30\) & Recorded value \\
\midrule
Producer word visits & \(143\,137\,713\) \\
Checker states summed over all layers & \(78\,459\,676\) \\
Checker transitions (state--label pairs) & \(84\,536\,537\) \\
Largest checker frontier (records) & \(11\,416\,082\) \\
Producer peak resident set (bytes) & \(6\,717\,906\,944\) \\
Checker peak resident set (bytes) & \(977\,747\,968\) \\
Logical coefficient table (bytes) & \(8\,589\,934\,592\) \\
\bottomrule
\end{tabular}
\caption{Work counters and memory quantities for the completed computation.
Checker counts come from the terminal checkpoint. The transition count
includes one contribution for each legal extension of a state by a label,
even when different extensions reach the same destination.}
\label{tab:work}
\end{table}

In Table~\ref{tab:work}, \(2^{29}\) cells of 16 bytes give an \(8\)-GiB
logical table. Peak resident memory was \(6.26\) GiB for the producer
and \(0.91\) GiB for the checker. The producer's zero-initialized allocation
can have pages that remain untouched during traversal, so its allocated
size need not all be resident. The coefficient check scans the full file once
per cardinality, requesting \(240\) GiB of logical reads at \(n=30\).
The ordered reduction rereads each relevant prefix in paired blocks.
These read volumes describe the algorithm's file accesses. Caching,
sorting and checkpoint I/O also affect the observed runtime.

\section{Pattern containment and lower-order kernels}

Hardin's original prefix gives \(a(n)=0\) for \(4\leq n\leq16\) and
\(a(17)=12\) \cite{hardin215}. The direct check covers this zero block.
The following realization exhibits positivity at \(n=17\):
\begin{align*}
 x&=(3,0,0,0,0,0,1,0,1,0,2,0,0,2,1,0,0,0,1,1,5),\\
 p&=(3,0,1,4,7,8,9,12,13,10,11,16,14,6,2,5,15).
\end{align*}
Substitution in \eqref{eq:definition} gives all 17 labels.
Standardizing any four entries of this permutation to their relative
ranks gives a length-four pattern. Heredity would require that pattern
to be realizable, whereas \(a(4)=0\).
Every count in the computed block \(17\le n\le30\) is positive,
but the counts can decrease within it, as \(a(20)=444>a(21)=420\) shows.
Eventual positivity and the complete set of zeros remain open here.

On a vector of length \(L\ge r+1\), let \(T_r\) form the \(L-r\) windows with
kernel \((\binom{r}{0},\ldots,\binom{r}{r})\).
Let \(\mathcal P_r(n)\) be the permutations of \(0,\ldots,n-1\) that
are images under \(T_r\) of nonnegative integer vectors of length \(n+r\).
The counts at orders \(2,3,4\) are A180213, A180214 and A180215,
respectively \cite{hardin213,hardin214,hardin215}.
If \(p=T_s x\) for \(s>r\), put \(z=T_{s-r}x\), of length \(n+r\).
Its entries are nonnegative integers. In the \(i\)-th window of \(T_rz\),
the coefficient of \(x_{i+j}\) is the coefficient of \(t^j\) in
\((1+t)^r(1+t)^{s-r}=(1+t)^s\), so \(T_rz=p\).
Thus \(\mathcal P_4(n)\subseteq\mathcal P_3(n)\subseteq\mathcal P_2(n)\)
and
\[
 a(n)\le \mathrm{A180214}(n)\le \mathrm{A180213}(n).
\]
Both inequalities are strict at \(n=4\), with counts \(0,2,6\) in
the respective OEIS entries \cite{hardin215,hardin214,hardin213}.
The inclusion uses \(z\) as an input to \(T_r\), and its entries
may repeat. For instance, \(x=(0,0,0,0,0,1)\)
gives \(p=(0,1)\) with quadratic intermediate vector \((0,0,0,1)\).
The lower-order counts restrict their outputs to permutations, so this
operator factorization does not give a composition rule for the counts.

\section{Conclusion}

The forced zero window reduces the count to a single table of one-sided
orderings, with multiplicities checked by the state recurrence.
This gives six further exact terms in the local \(b\)-file, ending at
\(a(30)=3209944\). The binomial kernel also orders the admissible sets
by inclusion across the three window orders.

\section*{Data availability}

The source and build instructions, direct-check program, canonical
\(b\)-file, result and validation tables, execution records, and manuscript
are available from the GitHub repository and Zenodo release series cited
in \cite{corti}. The Zenodo concept DOI identifies the series; individual
version records identify the corresponding immutable deposits.
The original coefficient table and traversal journal are retained locally
in a lossless archive of approximately \(224\) MiB
(\(8.53\) GiB expanded) for optional audits of the original execution.
They are excluded from the public package and can be regenerated from
the distributed sources. Both original archive members were verified in
full, and restoration was tested on small fixtures. Native audits of the
original run require the restored members and matching source, build,
executable and runtime identities. Source-based reproduction instead
creates fresh executables and scientific state from the defining inputs.

\section*{AI use disclosure}

Large language models assisted research and drafting for A180215.
OpenAI Codex supported the work. The code and the manuscript were
reviewed in separate rounds, with reported contributions across those
rounds from ChatGPT, Claude, Gemini, GLM, Grok, Meta AI, MiniMax and Qwen.
Tool names come from the project records, which did not consistently
identify exact model versions.

\smallskip
\noindent\textbf{Mathematical analysis.}
Codex assisted with deriving and checking the zero-window decomposition,
the side recurrence and the necessary arithmetic constraints.
NotebookLM was used to consult earlier projects on zero anchoring,
state multiplicities and checks from the original definition.
The claims are supported by the proofs and checked against the source
definition and historical OEIS data.

\smallskip
\noindent\textbf{Algorithm and code.}
Codex assisted with implementation and validation design.
The bounded direct C check was adapted from a MiniMax-supplied program
and executed during local code-review integration. Corti ran the frozen
producer and C++ checker.
The numerical claims rest on the local result tables, canonical
\(b\)-file, validator and execution records. The validators do not use
conversational assertions.

\smallskip
\noindent\textbf{Manuscript.}
Codex assisted with drafting, checking the exposition against the sources
and adjudicating the external manuscript reviews. The author takes full scientific
responsibility for the mathematical statements, algorithm, implementation,
numerical results, validation boundaries and conclusions.

\clearpage
\begin{thebibliography}{9}
\sloppy\raggedright
\bibitem{subsetconv}
A.~Bj\"orklund, T.~Husfeldt, P.~Kaski, and M.~Koivisto,
Fourier meets M\"obius: fast subset convolution,
in \emph{Proceedings of the 39th Annual ACM Symposium on Theory of Computing},
ACM, 2007, pp.~67--74.
\href{https://www.cs.helsinki.fi/u/mkhkoivi/publications/stoc-2007.pdf}{Author-hosted conference version}.
Preprint DOI \url{https://doi.org/10.48550/arXiv.cs/0611101}.

\bibitem{corti}
C. Corti, \emph{Permutations from Fourth-Order Binomial Window Sums:
OEIS A180215 through \(n=30\)}, GitHub repository and Zenodo archive, 2026.
\url{https://github.com/carcorti/A180215}\ and
\url{https://doi.org/10.5281/zenodo.24424435}.

\bibitem{gcc128}
Free Software Foundation, 128-bit Integers,
\emph{Using the GNU Compiler Collection (GCC)}, version 13.3.0.
\url{https://gcc.gnu.org/onlinedocs/gcc-13.3.0/gcc/_005f_005fint128.html}.

\bibitem{hardin213}
R.~H. Hardin, Sequence A180213, \emph{The On-Line Encyclopedia of
Integer Sequences}, 15 August 2010.
\url{https://oeis.org/A180213}.

\bibitem{hardin214}
R.~H. Hardin, Sequence A180214, \emph{The On-Line Encyclopedia of
Integer Sequences}, 15 August 2010.
\url{https://oeis.org/A180214}.

\bibitem{hardin215}
R.~H. Hardin, Sequence A180215, \emph{The On-Line Encyclopedia of
Integer Sequences}, 15 August 2010.
\url{https://oeis.org/A180215}.
\href{https://oeis.org/history?seq=A180215}{Entry history}.
\href{https://oeis.org/A180215/b180215.txt}{historical table}
(accessed 21 November 2026).
\end{thebibliography}

\end{document}
